% sections/bibliotheques.tex
% Phase 9: Bibliotheques, Outils et Technologies (RAPP-04)

\section{Biblioth\`eques, Outils et technologies}
\label{sec:bibliotheques}

Les biblioth\`eques et technologies utilis\'ees sont d\'ecrites ci-dessous,
organis\'ees par cat\'egorie.

\subsection{Analyse de graphe}

\paragraph{NetworkX (v3.4.2).}
NetworkX~\cite{networkx_docs} fournit les structures \texttt{DiGraph} et
\texttt{Graph} utilis\'ees dans ce projet, ainsi que les algorithmes de
centralit\'e (degr\'e, betweenness, closeness, PageRank), de d\'etection de
communaut\'es (Louvain, int\'egr\'e depuis la version~3.2), de composantes
connexes, de clustering et de plus courts chemins. C'est le socle commun
aux six questions de recherche.

\subsection{Visualisation}

\paragraph{Matplotlib (v3.9.4).}
Matplotlib~\cite{matplotlib_docs} g\'en\`ere les 16~figures PNG \`a 200~DPI.
Deux contextes graphiques sont utilis\'es~: fond sombre (\texttt{\#0F172A})
pour les graphes de r\'eseau, fond clair pour les graphiques analytiques.
La palette de couleurs (bleu \texttt{\#2563EB}, violet \texttt{\#7C3AED},
orange \texttt{\#F59E0B}) est centralis\'ee dans \texttt{src/style.py}.

\subsection{Traitement de donn\'ees}

\paragraph{Pandas (v2.2.3).}
Pandas charge et manipule le CSV de 401\,709 transactions~: conversion
des timestamps Unix en \texttt{datetime}, agr\'egation des ar\^etes multiples
via \texttt{groupby}, et fen\^etrage horaire pour l'analyse temporelle (RQ4).

\paragraph{NumPy (v2.1).}
NumPy fournit les tableaux et op\'erations vectoris\'ees pour les calculs
num\'eriques~: distributions de degr\'es, courbes de Lorenz, statistiques
descriptives (moyennes, m\'edianes, percentiles).

\paragraph{SciPy (v1.14).}
SciPy compl\`ete NumPy pour les tests statistiques, les comparaisons de
distributions et les corr\'elations de rang (Spearman).

\paragraph{powerlaw (v2.0.0).}
La biblioth\`eque \texttt{powerlaw} ajuste des distributions \`a queue
lourde~: estimation de l'exposant $\alpha$ et du seuil $x_{\min}$ par
maximum de vraisemblance, et tests de comparaison entre mod\`eles (loi de
puissance, lognormale, exponentielle), selon la m\'ethodologie de Clauset
et al.~\cite{clauset2009powerlaw}. Utilis\'ee en RQ1 et RQ6.

\subsection{R\'edaction et compilation}

\paragraph{\LaTeX{} avec pdf\LaTeX{}.}
Le rapport est compos\'e avec le syst\`eme \LaTeX{}, compil\'e par le moteur
pdf\LaTeX{}. Le document utilise la classe \texttt{article} avec 18~paquets
organis\'es selon un ordre de chargement strict~: \texttt{babel} pour les
conventions typographiques fran\c{c}aises, \texttt{amsmath} pour la notation
math\'ematique, \texttt{booktabs} pour les tableaux professionnels,
\texttt{graphicx} pour l'inclusion des figures, \texttt{biblatex} pour la
gestion bibliographique, \texttt{hyperref} pour les liens cliquables et
\texttt{cleveref} pour les r\'ef\'erences crois\'ees intelligentes. La
bibliographie est g\'er\'ee par \texttt{biblatex} avec le backend Biber,
offrant un support Unicode complet pour les noms d'auteurs comportant des
caract\`eres accentu\'es.
